\paragraph{Theorem F.2.39 (Proved: the full effective classification and all marked gluings).}
The linear classification and all-twist topological computation
asked for in F.2 have the following complete effective answer.
For $k=2$ and $k=3$, every admissible ordered integral tuple
has the canonical representative of F.2.37, with an explicit
simultaneous unimodular conjugator.  Equality of these
representatives is necessary and sufficient for simultaneous
integral conjugacy.  There is no fixed global height cutoff;
the finite stopping bound for an input of height $H$ is
$(2H+1)^{36(k+1)}$ candidates before the proved rejection
filters.  The separate involution normal form has exactly
56 classes and is computed completely in F.2.38.

For every choice of actual free affine orders and integer
numerators, integral clutch vector, and specified primitive
Kuhn cusp marking, F.2.31 and F.2.35 compute the full
fundamental group and every integral homology group of the
gluing.  This is the fully marked correction to F.2.3;
peripheral matrices are not substituted for a choice of
filling.  The algorithms have no cutoff on affine orders,
twist entries or lattice indices.  The complete list of data
with $\pi_1=1$ and $H_2=H_3=0$ is the constructive
parametrization of F.2.30, with its exact lattice enumeration
and necessary-and-sufficient finite tests.  Each retained
filling has the homology and oriented smooth type stated
there.

\begin{proof}
The input is finite, so the prime-subgroup test in F.2.31
decides whether every claimed affine cyclic action is free
and has the stated actual order.  The integral cusp test and
its adapted basis are F.2.1.  The complete linear admissibility
filter and the canonical representative, including the proved
finite stopping bound, are F.2.2 and F.2.36--38.  The
characteristic-polynomial filters used by that algorithm are
conjugacy invariants, so none can remove a possible
representative.  The unbounded trace family (F.CN3) proves
that no fixed finite box could replace the input-dependent
bound.

We spell out how the linear classification is used with a
specified gluing.  If $R$ is its returned conjugator and
$P$ is the supplied cusp marking, change the data by
\[
 M_i'=RM_iR^{-1},\qquad v_i'=Rv_i,\qquad
 \kappa'=R\kappa,\qquad P'=RP.                           \tag{F.SY1}
\]
The actual orders are unchanged.  The coordinate map
$x\mapsto Rx$ on the torus intertwines the finite quotient
charts, since it sends $M_ix+v_i/m_i$ to
$M_i'x'+v_i'/m_i$.  At the cusp, its coordinates in the
marked Kuhn model are unchanged:
$(P')^{-1}x'=P^{-1}x$.  Hence it intertwines every toric
chart, translated face identification and collapse map of
that fixed filling.  It also carries the clutch path to
the path with vector $\kappa'$.  These maps agree on the
collars and give a diffeomorphism of the two fully marked
gluings.  Thus homology can be computed in the original
coordinates, while the canonical data and their group
presentation are also returned.  Independently replacing
$P'$ by another Smith marking is not this comparison.

An actual order-one site is ordinary, even if its integer
numerator has been recorded.  Delete it by the explicit
comparison of F.2.31: if it is the $i$th site, replace
$\kappa$ by
$\kappa+(M_{i+1}\cdots M_k)^{-1}v_i$ before deletion.
The cusp matrix and its supplied marking are unchanged.
There cannot be an order-one site in an admissible two-point
tuple: its other finite matrix would be the inverse of a
nontrivial unipotent matrix.  Nor can a three-point tuple
have two such sites.  Thus this reduction leaves precisely
the genuine two- or three-point inputs of F.2.35 and loses
no case from the prompt.

For each remaining input, the exact unsaturated normal
closure, cyclic or triangle-group extension and exceptional
spherical Euler relation of F.2.31 give the full fundamental
group.  The finite cap, complement and marked Kuhn chain
maps of F.2.33--35 give a finite integer Mayer--Vietoris cone
through degree nine.  Integer kernels and Smith reduction
therefore give all $H_0,\ldots,H_8$ with their complete
torsion factors.  This is a calculation for arbitrary
supplied twists, rather than a restriction to the affine
orders of one displayed example.  Its independent
abelianization check compares the resulting $H_1$ with the
van Kampen answer.

Finally F.2.7--30 prove both directions of the survivor
classification.  Genuine three-point orbifold quotients
exclude simple connectivity.  The two-point filters reduce
every survivor to the minimal order-$(5,4)$ spectrum; the
rational normal form, finite stable-lattice and extension
algorithms include every integral representation above it.
The universal affine normalization includes every twist,
and the six exact integral tests decide retention.
Conversely every retained datum satisfies all three
topological conditions by those tests.  F.2.30 computes
its common integral homology and applies the almost-complex
and smooth-classification arguments to identify its oriented
type.  There is no residual untested rational-height bound
in this parametrization.  The canonical-representative
algorithm can remove duplicates from it if desired, but
no finite global enumeration is asserted or needed for
these complete per-input decisions.
\end{proof}

The unified entry point is
\path{code/fourfolds_classify_gluing.py}.
It records the original and transported markings, the exact
linear normalizer, both group presentations and the complete
integer homology calculation.  Its general noninvolution
canonical branch has the proved finite stopping rule of
F.2.37 and the same stated practical limitation: that
astronomically large fallback is not claimed exhaustively
executed.  The executed finite catalogue, conjugacy tests,
full marked homology examples and their exact bounds are
those recorded in F.2.33--38.
The identity-insertion branch has additionally been run on
all 168 inputs obtained by placing an identity in each of
the three marked positions of each of the 56 involution
classes, after a nontrivial simultaneous conjugation.
Every exact representative and unimodular conjugator passes.
The full unified driver has been run on the orientation-reversing
involution gluing of F.2.35 and on the same gluing presented
with an additional actual order-one site.  In the latter
test its numerator is $(0,1,0,1,0,0)^T$ and the original
clutch is $(1,-1,0,1,0,0)^T$; the suffix comparison reduces
the clutch to zero.  Both executions return the complete
canonical data, transported cusp marking, exact group and
the full integral homology of that example, with the
independent abelianization check.  Their outputs and full
JSON certificates are retained.
